% Shared preamble for "Mathematics for the Observation Programme".
% Every chapter uses only the packages and macros defined here.
\usepackage[T1]{fontenc}
\usepackage[utf8]{inputenc}
\usepackage{lmodern}
\usepackage{microtype}
\usepackage[margin=1.1in]{geometry}
\usepackage{amsmath,amssymb,amsthm,mathtools}
\usepackage{bm}
\usepackage{booktabs}
\usepackage{array}
\usepackage{enumitem}
\usepackage{xcolor}
\usepackage{graphicx}
\usepackage{tikz}
\usepackage{pgfplots}
\pgfplotsset{compat=1.18}
\usetikzlibrary{arrows.meta,positioning,calc,decorations.pathreplacing,patterns,shapes.geometric,fit,backgrounds}
\usepgfplotslibrary{fillbetween,groupplots}

% ---------- figure palette and styles (use these in every figure) ----------
\definecolor{pblue}{RGB}{31,94,166}     % primary: the source, the signal, the main curve
\definecolor{porange}{RGB}{214,110,22}  % the consumer, the read operator, the reader
\definecolor{pgreen}{RGB}{46,139,87}    % the optimum, the accepted, the kept
\definecolor{pred}{RGB}{192,40,45}      % the error, the floor, the rejected
\definecolor{pgray}{RGB}{120,120,120}   % axes, guides, context
\definecolor{plight}{RGB}{232,238,247}  % fills and shaded regions
\tikzset{
  vec/.style={-{Stealth[length=2.6mm]},thick},
  guide/.style={pgray,dashed,thin},
  block/.style={draw,rounded corners=2pt,thick,fill=plight,minimum height=7mm,minimum width=16mm,align=center,font=\small},
  flow/.style={-{Stealth[length=2.2mm]},thick},
  world/.style={circle,draw,thick,minimum size=7mm,inner sep=0pt,font=\small},
  lbl/.style={font=\small},
}
\pgfplotsset{
  primer/.style={
    width=0.62\linewidth, height=0.40\linewidth,
    axis lines=left, tick label style={font=\footnotesize},
    label style={font=\small}, legend style={font=\footnotesize,draw=none,fill=none},
    every axis plot/.append style={thick},
    cycle list={{pblue},{porange},{pgreen},{pred},{pgray,dashed}},
  },
}
\usepackage[numbers,sort&compress]{natbib}
\usepackage[colorlinks=true,linkcolor=blue!50!black,citecolor=green!35!black,urlcolor=blue!50!black]{hyperref}
\usepackage[nameinlink,capitalise]{cleveref}

\numberwithin{equation}{chapter}

% ---------- environments ----------
% Definitions, examples and results share one numbering sequence. Each gets
% an alias counter so that \cref prints its own name instead of "Definition".
\usepackage{aliascnt}
\newcommand{\sharedtheorem}[3]{% #1 env, #2 name, #3 plural
  \newaliascnt{#1}{definition}%
  \newtheorem{#1}[#1]{#2}%
  \aliascntresetthe{#1}%
  \crefname{#1}{#2}{#3}\Crefname{#1}{#2}{#3}}
\theoremstyle{definition}
\newtheorem{definition}{Definition}[chapter]
\crefname{definition}{Definition}{Definitions}
\sharedtheorem{example}{Example}{Examples}
\newtheorem{exercise}{Exercise}[chapter]
\crefname{exercise}{Exercise}{Exercises}
\theoremstyle{plain}
\sharedtheorem{theorem}{Theorem}{Theorems}
\sharedtheorem{proposition}{Proposition}{Propositions}
\sharedtheorem{lemma}{Lemma}{Lemmas}
\sharedtheorem{corollary}{Corollary}{Corollaries}

% Running heads. Chapter and section titles in this book are full
% sentences, too long for a page header, so the header shows numbers only.
\renewcommand{\chaptermark}[1]{\markboth{\MakeUppercase{\chaptername\ \thechapter}}{}}
\renewcommand{\sectionmark}[1]{\markright{\MakeUppercase{Section\ \thesection}}}
\theoremstyle{remark}
\newtheorem*{remark}{Remark}

% "Builds on" and "Used in" pointers at the end of a section
\newcommand{\buildson}[1]{\par\noindent\textit{Builds on.} #1\par}
\newcommand{\usedin}[1]{\par\noindent\textit{Used in.} #1\par}
% A boxed "where this appears in the projects" pointer
\newcommand{\inproject}[1]{\par\smallskip\noindent\fbox{\parbox{0.96\linewidth}{\small\textbf{In the projects.} #1}}\par\smallskip}

% ---------- notation (see appendix A) ----------
\newcommand{\R}{\mathbb{R}}
\newcommand{\N}{\mathbb{N}}
\newcommand{\E}{\mathbb{E}}
\newcommand{\Prob}{\Pr}
\DeclareMathOperator{\tr}{tr}
\DeclareMathOperator{\rank}{rank}
\DeclareMathOperator{\range}{range}
\DeclareMathOperator{\diag}{diag}
\DeclareMathOperator{\Var}{Var}
\DeclareMathOperator{\Cov}{Cov}
\DeclareMathOperator*{\argmin}{arg\,min}
\DeclareMathOperator*{\argmax}{arg\,max}
\renewcommand{\ker}{\operatorname{ker}}
\newcommand{\T}{^{\top}}
\newcommand{\inner}[2]{\langle #1,#2\rangle}
\newcommand{\norm}[1]{\lVert #1\rVert}
\newcommand{\abs}[1]{\lvert #1\rvert}
\newcommand{\psd}{\succeq}          % Loewner order
\newcommand{\pd}{\succ}
\newcommand{\Sym}{\mathbb{S}}       % symmetric matrices, \Sym^d_+ PSD cone
\newcommand{\PC}{P_C}               % read operator of consumer C
\newcommand{\Sd}{\Sigma_\delta}     % error covariance
\newcommand{\dO}{d_O}               % observer distortion tr(P_C Sigma_delta)
\newcommand{\RC}{R_C}               % consumer-relative rate-distortion function
\newcommand{\KL}[2]{D\!\left(#1\,\middle\|\,#2\right)}
\newcommand{\MI}[2]{I(#1;#2)}
\newcommand{\Ent}{H}
\newcommand{\h}{h}                  % differential entropy
\newcommand{\Normal}{\mathcal{N}}
\newcommand{\iid}{\overset{\text{iid}}{\sim}}
